\subsection{Modelos de t\'erminos}
Sea $\lang := (\lambda, \mu, K)$ un lenguaje formal, definimos
\[ \kappa_\lang = \max \{ \aleph_0, |I|, |J|, |K| \}. \]
La cardinalidad del lenguaje $\lang$, mientras que el cardinal de una $\lang$-estructura es el cardinal del universo.
\footnote{Si consideramos un lenguaje a la Ziegler $\kappa_\lang := \max \{ \lang, \aleph_0 \}$.}
Tenemos que 
\begin{align*}
  |\mathrm{Tm} (\lang)| &= \max \{ \aleph_0, |J|, |K| \},\\
  |\mathrm{Fml} (\lang)| &= \kappa_\lang,\\
  |\mathrm{Sent} (\lang)| &= \kappa_\lang.
\end{align*}
Sea $\Sigma \subseteq \mathrm{Sent} (\lang)$ consistente, tenemos que la extensi\'on $\lang'$ de $\lang$ tiene el mismo cardinal.
Trivialmente
\[ \kappa_\lang \leq \kappa_{\lang'}. \]
Luego, procedemos por inducci\'on, supongamos que 
\[ \kappa_{\lang_n} \leq \kappa_\lang, \]
de esto se sigue que 
\[ \kappa_{\lang_n+1} \leq \max \{ \kappa_{\lang_n}, M_n \} \leq \kappa_{\lang_n} \leq \kappa_\lang.  \]
Por lo tanto $\kappa_{\lang_n} \leq \kappa_\lang$ para todo $n \in \nn$.
De esto se sigue que $\kappa_{\lang'} \leq \kappa_\lang$.
Por lo tanto $|\mathrm{CT} / \approx| \leq \kappa_\lang$.
\begin{theorem}
  Sea $\Sigma \subseteq \mathrm{Sent} (\lang)$ tal que tiene un modelo de cardinal infinito y sea $\kappa$ un cardinal tal que $\kappa_\lang \leq \kappa$.
  Existe un modelo de $\Sigma$ con cardinal $\kappa$. 
\end{theorem}\\
\emph{Demostraci\'on.}\\
Definamos $\lang_\kappa$ como $I_\kappa := I, J_\kappa := J$ y $K_\kappa := K \cup \kappa$.
\footnote{Viendo a $\kappa$ como ordinal.}
Extendemos $\Sigma$ a $\Sigma_\kappa := \Sigma \cup \{ c_\nu \neq c_\mu \ | \ \nu < \mu < \kappa \}$. 

Sea $\Pi$ un subconjunto finito de $\Sigma_\kappa$, como $\Sigma$ tiene un modelo infinito, este mismo modelo lo podemos extender a un modelo de $\Pi$. 
Por compacidad, $\Sigma_\kappa$ es consistente.
El modelo de t\'erminos
\footnote{El usado en el teorema de completud de G\"odel}
tiene cardinal menor o igual a $\kappa$ y como tiene $\kappa$ t\'erminos, su cardinal es $\kappa$. \qed
\newpage
\begin{proposition}
  Sean $\mathfrak{A}$ y $\mathfrak{B}$ $\lang$-estructuras si $\mathfrak{B} \models \mathrm{Th} (\mathfrak{A})$, entonces $\mathfrak{A}$ y $\mathfrak{B}$ son esencialmente equivalentes.
\end{proposition} \hfill \break
\emph{Demostraci\'on.} \hfill \break
Por hip\'otesis tenemos que para toda $\varphi \in \mathrm{Sent} (\lang)$, 
\[ \mathfrak{A} \models \varphi \Longrightarrow \mathfrak{B} \models \varphi. \]
Luego, supongamos que existe $\psi \in \mathrm{Sent} (\lang)$ tal $\mathfrak{A} \not\models \psi$, es decir, $\mathfrak{A} \models \neg \psi$.
Por lo tanto $\mathfrak{B} \models \neg \psi$ y por definici\'on $\mathfrak{B} \not\models \psi$.
Concluyendo que 
\[ \mathfrak{A} \not \models \psi \Longrightarrow \mathfrak{B} \not\models \psi. \]
Lo cual demuestra por contrapositiva la segunda implicaci\'on. \qed
\begin{corollary}
  Sea $\mathfrak{A}$ una $\lang$-estructura infinita y $\kappa \geq \kappa_\lang$ un cardinal, existe una $\lang$-estructura $\mathfrak{B}$ de cardinal $\kappa$ tal que es elementalmente equivalente a $\mathfrak{A}$.
\end{corollary} \hfill \break
\emph{Demostraci\'on.} \hfill \break
Se sigue del teorema anterior cuando consideramos a $\Sigma = \mathrm{Th} (\mathfrak{A})$, tenemos que existe un modelo de $\mathrm{Th} (\mathfrak{A})$ de cardinal $\kappa$ y de la proposici\'on anterior llegamos a que es elementalmente equivalente a $\mathfrak{A}$. \qed

